\subsection{Covariance matrix}
A major computational bottleneck of of the real-time streaming algorithm is the recalculation of the covariance matrix at each time step as the history window advances.
To solve this issue, the implementation exploits the overlap of a sliding window: When a new value is appended to the back of the window, a single embedding vector, named $tde_{old}$ is dropped from the active window and a new vector $tde_{new}$ is formed. All intermediate time delay embedded vectors within the window stay exactly the same, their respective structural contributions to the total covariance does thereby also not change. 
The contribution of a single multidimensional vector to the covariance matrix is found using an outer product $x{x}^T$. The incoming time delay embedding vectors all have the same dimension $d$ multiplying it by its own transpose it creates an $d*d$ square matrix. \\
\begin{center}
    ${tde} \cdot {tde}^T = 
\begin{bmatrix} 
x \\ 
y \\ 
z 
\end{bmatrix}
\begin{bmatrix} 
x & y & z 
\end{bmatrix} 
= 
\begin{bmatrix} 
x^2 & xy & xz \\ 
yx & y^2 & yz \\ 
zx & zy & z^2 
\end{bmatrix}$
\end{center}

Instead of calculating the covariance matrix from scratch each time new data arrives we introduce a running $d*d$ matrix which is kept and updated along each iteration. As new vectors arrives the running covariance matrix does only have to update the incoming and outgoing dot products. 
The running covariance matrix is maintained by incrementally updating the sum of outer products (scatter matrix) $S_t$. When the sliding window reaches a steady state, the matrix is updated by subtracting the outer product of the dismissed, centered vector $\mathbf{x}_{\text{old}}$ and adding the outer product of the incoming centered vector $\mathbf{x}_{\text{new}}$:
\begin{equation}
    \mathbf{S}_t = \mathbf{S}_{t-1} - \left( \mathbf{x}_{\text{old}} \mathbf{x}_{\text{old}}^T \right) + \left( \mathbf{x}_{\text{new}} \mathbf{x}_{\text{new}}^T \right)
    \label{eq:running_covariance_update}
\end{equation}
where $x_{new}, x_{old}$ represent the zero-mean centered time-delay embedding vectors. 
The unbiased sample covariance matrix $\mathbf{\Sigma}_t$ is subsequently obtained by scaling the scatter matrix by the number of active degrees of freedom:
\begin{equation}
    \mathbf{\Sigma}_t = \frac{1}{N - 1} \mathbf{S}_t
    \label{eq:covariance_scaling}
\end{equation}
where $N$ denotes the current sample count in the window ($N \ge 2$).his normalization ensures that the computed eigenvalues reflect true data variance independently of the active window sample size $N$.
With this modification the complexity isn't dependent of the window size $W$ anymore but only the time delay embedding dimensions $d$.This improves the complexity from $O(N*(d*d))$ to $O(d*d)$ which means that the window can become really large while keeping the speed constant to the dimensions used. 